\section{Results}

\subsection{Identity adjudication increases the reference scope}

The raw large-lake scope contained 179{,}787 HydroLAKES records. Of the 1{,}998 records removed by the upstream proximity screen, 108 had authoritative duplicate evidence and 645 satisfied the calibrated spatial-and-area rule. Restoring the other 1{,}245 records produced a reference inventory of 199{,}976 water bodies (Fig.~\ref{fig:framework}a). Retaining the 645 spatial matches in an authoritative-only sensitivity increased the inventory to 200{,}621 but added only 14.9~TWh~yr$^{-1}$ to L0 and 14.7~TWh~yr$^{-1}$ to L1. Identity policy therefore had a small energy effect but a material effect on inventory traceability.

Under the reference coverage scenario, L0 was 16{,}713~TWh~yr$^{-1}$ (Fig.~\ref{fig:framework}b). L1 retained 79{,}669 water bodies and 10{,}995~TWh~yr$^{-1}$, or 65.79\% of L0. Reservoirs retained 5{,}194 of 5{,}335~TWh~yr$^{-1}$ (97.37\%), whereas natural and controlled lakes retained 5{,}800 of 11{,}378~TWh~yr$^{-1}$ (50.98\%). The unequal retention moved the climate-screened resource away from the high-latitude lake concentrations that dominate the unscreened inventory.

\subsection{Climate screening changes the regional ordering}

North America contained 7{,}607~TWh~yr$^{-1}$ at L0 but retained 3{,}354~TWh~yr$^{-1}$ at L1 (Fig.~\ref{fig:geography}a,b). Europe contracted from 2{,}543 to 1{,}307~TWh~yr$^{-1}$. Asia retained 3{,}112 of 3{,}325~TWh~yr$^{-1}$, while South America retained 1{,}467 of 1{,}483~TWh~yr$^{-1}$. Africa and Oceania showed no material L1 loss under the reference hierarchy.

Canada led L0 with 5{,}216~TWh~yr$^{-1}$ but fell to 1{,}210~TWh~yr$^{-1}$ at L1. The United States then ranked first at 1{,}797~TWh~yr$^{-1}$, followed by Canada, China (1{,}190~TWh~yr$^{-1}$), Brazil (693~TWh~yr$^{-1}$) and Australia (606~TWh~yr$^{-1}$) (Fig.~\ref{fig:geography}c). Thus, inventory abundance and climate-screened opportunity produced different geographic priorities.

\subsection{Coverage and ice policy dominate the tested resource sensitivities}

The specific-yield benchmark contained 134{,}703 matched lakes. Annual yield correlated with the published Woolway output at $r=0.896$, with a mean bias of $-144$~kWh~kWp$^{-1}$, mean absolute error of 159~kWh~kWp$^{-1}$ and median relative difference of $-11.89\%$ (Fig.~\ref{fig:validation}a). Scaling POWER wind to the sampled Google Earth Engine mean changed L0 by $-0.70\%$ on the common legacy climate scope; a zero-wind stress test changed it by $-1.39\%$.

Coverage assumptions had a much larger effect (Fig.~\ref{fig:validation}b). Uniform 5\%, 10\%, 20\% and 30\% coverage produced L0 values of 7{,}743, 13{,}909, 24{,}606 and 34{,}190~TWh~yr$^{-1}$, respectively. The type-specific reference scenario produced 16{,}713~TWh~yr$^{-1}$. Its L1 retention (65.79\%) was also scenario-dependent because the 30~km$^2$ cap changes the relative weight of large cold lakes.

The reference ice hierarchy retained 10{,}995~TWh~yr$^{-1}$ (Fig.~\ref{fig:validation}c). Removing LI-CCR priority and using Woolway followed by the temperature proxy retained 10{,}848~TWh~yr$^{-1}$, 1.33\% less. Applying the temperature proxy to all records retained 12{,}562~TWh~yr$^{-1}$, 14.26\% more. LI-CCR and Woolway classifications disagreed for 2{,}893 of 22{,}775 records with both sources. The hierarchy therefore affects the L1 endpoint even though the reference and Woolway-plus-proxy totals are close.

\subsection{Public surface-area evidence defines L2 intervals}

After L1, the core conservation screen retained 68{,}111 water bodies and 8{,}182~TWh~yr$^{-1}$; the population rule reduced this to 33{,}157 water bodies and 5{,}058~TWh~yr$^{-1}$ (Fig.~\ref{fig:l2_evidence}a). Within that set, 22{,}405 records had complete 1991--2018 evidence and no observed zero-area month, contributing 4{,}190~TWh~yr$^{-1}$. A further 4{,}687 records (254~TWh~yr$^{-1}$) had an observed zero-area month and failed both endpoints. The 6{,}065 unknown records contributed 614~TWh~yr$^{-1}$. The core L2 interval was consequently 4{,}190--4{,}805~TWh~yr$^{-1}$, or 25.07--28.75\% of L0 (Fig.~\ref{fig:l2_evidence}b).

The conservative conservation screen retained 56{,}865 water bodies and 6{,}047~TWh~yr$^{-1}$; the population rule retained 27{,}840 and 3{,}607~TWh~yr$^{-1}$. Complete, no-dry-up evidence contributed 2{,}914~TWh~yr$^{-1}$, and 490~TWh~yr$^{-1}$ remained unknown. The conservative L2 interval was 2{,}914--3{,}404~TWh~yr$^{-1}$ (17.44--20.37\% of L0).

The interval width came entirely from the reservoir branch because the retained lake records had complete GLEV links. Core reservoir generation ranged from 2{,}594 to 3{,}208~TWh~yr$^{-1}$, whereas lake generation remained 1{,}596~TWh~yr$^{-1}$. Alternative dry-up criteria placed the core lower endpoint between 4{,}017 and 4{,}232~TWh~yr$^{-1}$ (Fig.~\ref{fig:l2_evidence}c). This criterion sensitivity does not resolve the universal 2019--2020 gap.

\subsection{Infrastructure mapping drives the partial-L3 range}

The core partial-L3 lower endpoint combined known no-dry-up evidence with conservative road-distance and observed power-line-distance policies. It retained 11{,}084 water bodies and 1{,}595~TWh~yr$^{-1}$ (9.54\% of L0). The upper endpoint combined unknown-pass evidence with possible-distance policies and retained 20{,}393 water bodies and 3{,}508~TWh~yr$^{-1}$ (20.99\%) (Fig.~\ref{fig:accessibility}a). The conservative interval was 1{,}101--2{,}494~TWh~yr$^{-1}$ for 8{,}549--16{,}489 water bodies (6.58--14.92\%).

The core factorial separated evidence and mapping policies (Fig.~\ref{fig:accessibility}b). With the evidence policy fixed at known-only, changing lower to upper mapping increased the result from 1{,}595 to 3{,}116~TWh~yr$^{-1}$. With lower mapping fixed, changing known-only to unknown-pass increased it from 1{,}595 to 1{,}786~TWh~yr$^{-1}$. The corresponding contrasts at the other factorial levels were 1{,}722 and 393~TWh~yr$^{-1}$. Mapping policy was therefore the larger source of the core partial-L3 spread.

GRanD hydropower-use evidence added 108~TWh~yr$^{-1}$ to the core lower-mapping endpoint and 304~TWh~yr$^{-1}$ to its upper-mapping counterpart (Fig.~\ref{fig:accessibility}c). Mapped lines alone retained 1{,}487 and 2{,}812~TWh~yr$^{-1}$ under the respective known-only cases. These quantities describe mapped infrastructure association and proximity; complete project-level L3 remains unavailable.

\section{Discussion}

The analysis changes the global FPV question from ``How large is the deployable potential?'' to ``Which part of a declared resource is supported for a specific screening decision?'' The reference sequence contracts from 16{,}713~TWh~yr$^{-1}$ at L0 to 65.79\% at L1, 17.44--28.75\% across the two L2 definitions, and 6.58--20.99\% after partial infrastructure screening. This contraction is large enough to reject the earlier manuscript structure in which successive levels differed only marginally. It also shows why a single endpoint cannot represent resource, suitability and deployability simultaneously.

Coverage remains the dominant design assumption before exclusions are applied. The tested uniform scenarios span more than a factor of four, consistent with the broad variation observed among operating FPV plants~\citep{nobre2024}. The 30\% reservoir and 10\% lake fractions are therefore useful as a transparent comparison case, not as a global recommendation. This distinction also clarifies comparisons with reservoir-centred estimates~\citep{lee2020,jin2023} and lake-centred estimates~\citep{woolway2024}: differences reflect scope, coverage and cap choices before they reflect model performance. Reporting the full coverage response is more informative for planners than presenting one fraction as technically optimal.

The L1 results reveal where inventory expansion changes the governing constraint. Reservoir generation was largely retained, whereas roughly half of lake generation failed the reference ice hierarchy. Canada and Russia consequently fell sharply after L1, while lower-latitude reservoir-rich countries rose in the screened ranking. The cross-model yield correlation supports the large-scale spatial pattern, but the negative yield bias and the disagreement between ice sources argue against site ranking by model output alone. Global screening can identify resource regions; hourly production, ice loading and system design still require local models and observations.

The L2 bounds turn missing surface-area evidence into an actionable data priority. Their width is concentrated in reservoirs without complete lake-linked histories rather than distributed uniformly. Improving reservoir-to-lake crosswalks and extending public monthly surface-area records would therefore narrow the screening interval directly. The lower endpoint is appropriate when a programme requires observed support for every global gate. The upper endpoint is useful when the purpose is to identify locations where new surface-area evidence could change a decision. Neither endpoint reproduces the published 1991--2020 criterion because the public GLEV evidence ends in 2018~\citep{zhao2022glev}.

The partial-L3 factorial provides a second planning lesson. Uncertainty in the spatial interpretation of road cells and sampled power lines was larger than the tested unknown-record policy. Higher-resolution roads, substations, voltage levels and interconnection queues would thus add more decision value than simply replacing an unknown with a favourable assumption. Hydropower association can prioritize follow-up at existing energy-water assets, but it does not demonstrate spare connection capacity. A practical use is to treat lower-endpoint sites as candidates for regional portfolio screening and upper-only sites as a data-acquisition queue.

Several global gaps remain outside the present hierarchy. GLOBathy maximum depth cannot resolve array or anchor conditions, and the parametric cost cases omit connection, anchoring, permitting and financing terms. Wind and wave extremes, water-level range, bathymetry, navigation, ecological response, water rights, social acceptance and project-specific grid studies are not globally complete~\citep{worldbank2019fpv}. These omissions are not additional numerical discounts that can be inferred from the available layers. They define the transition from global screening to project development. Future work should extend surface-area time series, publish auditable grid-capacity data, validate production against operating FPV fleets, and connect global candidates to site-scale engineering and environmental assessments.
